\section{Related Work and Current Field Position}
\label{sec:related_work}

\subsection{Peer-reviewed research systems}

The current field contains several peer-reviewed systems with wider research
scope than the controller studied here. The 2026 Nature version of The AI
Scientist reports automated ideation, experimentation, writing, and review in
machine learning~\cite{lu2026aiscientist}. One of three selected manuscripts
passed review at a workshop with a 70\% acceptance rate; the authors also
report implementation, rigor, hallucination, and citation failures, and do
not claim a main-conference standard. Co-Scientist uses asynchronous agents,
debate, ranking, evolution, memory, expert steering, and experimental
feedback for hypothesis development~\cite{gottweis2026coscientist}. Its role
is collaborative and retains expert evaluation. Robin integrates literature
search, hypothesis generation, human-selected candidates, human-run
experiments, and iterative analysis in experimental biology
~\cite{ghareeb2026robin}. ERA uses tree search to optimize empirical software
against quantitative objectives~\cite{aygun2026era}. Its authors distinguish
this objective from a claim of scientific discovery.

Earlier peer-reviewed work addresses narrower components. ResearchAgent and
SciMON use retrieved literature and iterative feedback for research-idea
generation~\cite{baek2025researchagent,wang2024scimon}. Box's Loop searches
over statistical models and analyses~\cite{li2024boxlm}. ReAct and Reflexion
provide reasoning-action and feedback-memory patterns used by many agent
controllers~\cite{yao2023react,shinn2023reflexion}. Against this background,
the present state machine is not positioned as a broader automation system.
Its object of study is the enforceable boundary between a negative result, a
successor mechanism, and a report claim.

\subsection{Independent evaluation and audit evidence}

AstaBench is a peer-reviewed ICLR 2026 benchmark suite covering research
planning, literature tasks, coding, and data analysis across 57 agents and 22
agent classes~\cite{bragg2026astabench}. Its results support component-level
progress but not a solved scientific-assistance problem. REPRO-Bench is a
peer-reviewed ACL 2025 study of 112 paper-and-package reproducibility
assessments; the strongest representative baseline reported 21.4\%
accuracy~\cite{hu2025reprobench}. Both studies separate a benchmark outcome
from a broad statement about discovery.

Several current sources are preprints and are treated here as provisional
evidence. The CORE-Bench preprint defines 270 reproducibility tasks from 90
papers and reports 21\% accuracy for its strongest baseline on the hardest
level~\cite{siegel2024corebench}. The PaperBench preprint evaluates replication
of 20 machine-learning papers and reports a 21\% average score for its
strongest tested agent, below its recruited human baseline
~\cite{starace2025paperbench}. An independent preprint audit of KOSMOS reports
a mixture of supported, uncertain, and false hypotheses in one radiation
biology assessment~\cite{nusrat2025kosmosaudit}. A 2026 preprint benchmark on
academic integrity studies whether research agents report failure honestly
under pressure~\cite{yang2026sciintegrity}. These sources motivate external
checks and adversarial failure tests; their review status prevents treating
their numerical findings as settled field estimates.

Evaluation design also matters. The AI Agents That Matter preprint argues for
attention to cost, model-scaffold attribution, meaningful downstream
outcomes, and reproducible comparisons~\cite{kapoor2024agentsmatter}. The
analysis by Demsar of comparisons across data sets supports treating the task
or data set, rather than repeated runs within it, as the unit for many
cross-task comparisons~\cite{demsar2006}. This is the statistical reason for
the additive task-level analysis in Section~\ref{sec:experimental_design}.

\subsection{Provenance and open research objects}

Evidence graphs need interoperable semantics in addition to local hashes.
The W3C PROV-O Recommendation defines entities, activities, agents, and their
provenance relations~\cite{lebo2013provo}. RO-Crate 1.3 provides a
machine-actionable research-object packaging specification
~\cite{rocrate2026spec}. Workflow Run RO-Crate profiles prospective and
retrospective workflow provenance and maps its terms to W3C PROV
~\cite{leo2024workflowrun}. The FAIR principles for research software add
software-specific expectations for findability, accessibility,
interoperability, and reuse~\cite{barker2022fair4rs}.

The current package is hash-linked and locally reproducible, but it is not
claimed to conform to those interoperable profiles. That mapping is separate
future work because adding metadata cannot convert an author rerun into
independent validation. ACM artifact policy likewise distinguishes artifact
availability and functionality from independently validated results
~\cite{acm2020artifacts}. This separation is retained throughout the paper.
